\documentclass[a4paper,11pt]{article}
\usepackage[ngerman]{babel}
\usepackage[top=3cm,bottom=3cm,left=3cm,right=3cm]{geometry} %NICHT IN KEMIE2


%\usepackage{microtype} %schöneres typesetting  mit [stretch=10] für nur 1% anstelle von 2%
%\usepackage{lmodern} %Latin Modern FONT
%\usepackage{pifont} %schönere Font



\usepackage{amsmath}
\usepackage{nicefrac}
\usepackage{amssymb}
%\usepackage{mathtools}

\usepackage{titlesec} %Um Section, etc bearbeiten zu können
\usepackage{setspace}
\usepackage{enumitem} %enumerate
\usepackage{multicol}
\usepackage[many]{tcolorbox}
\usepackage{hyperref}
\usepackage{arydshln}

%\renewcommand{\ttdefault}{lmtt} %Schriftart wird geändert zu Latin Modern Typewriter


\hypersetup{hidelinks,
			colorlinks=true,
			linkcolor=black,
			citecolor=miudiutex,
			urlcolor=miugray2
			}

\setlength{\parindent}{0pt}
\setcounter{section}{0}


\titlespacing*{\section}{0pt}{20mm}{6mm}
\titlespacing*{\subsection}{0pt}{6mm}{3mm}
\titlespacing*{\subsubsection}{0pt}{3mm}{0mm}



%\definecolor{miudiutex}{HTML}{2f3d39}
\definecolor{miudiutex}{HTML}{1d2926}
\definecolor{miugray}{HTML}{b4cccf}
\definecolor{miugray2}{HTML}{4d7365}
\definecolor{red}{HTML}{cf1f5c}
\definecolor{green}{HTML}{00A36C}
\definecolor{delphinidin}{HTML}{315794}
\definecolor{pelargonidin}{HTML}{b33807}
\definecolor{cyanidin}{HTML}{ba0746}
\definecolor{petunidin}{HTML}{b56fed}



\raggedcolumns %wichtig für Multicolumns


%================================
% SYMBOLS
%================================

\newcommand{\RA}{$\rightarrow$~}
\newcommand{\RL}{$\rightleftharpoons$~}
\newcommand{\HARP}{\ding{226}~}
%$\xrightarrow{2}$ %Pfeil mit 2 drüber


%================================
% SHORT COMMANDS (BOXES ; ETC)
%================================


\newcommand{\notiz}[1]{\begin{notizz}{}{}\begin{center}
			{\color{miugray2!50!miudiutex}#1} \end{center}\end{notizz}}
\newcommand{\zit}[2]{\begin{zitat*}{#1}{}\small #2 \end{zitat*}}

\newcommand{\frage}[2]{\begin{qst*}{#1}{}\small #2 \end{qst*}}

\newcommand{\unten}{\end{center} \tcblower \begin{center}}


\newcommand*\circled[1]{\tikz[baseline=(char.base)]{
            \node[shape=circle,draw,inner sep=0.5pt,miudiutex] (char) {\tiny #1};}}
            
\newcommand{\miusquare}{{\color{miugray2!50}\scriptsize $\blacksquare$}~}
\newcommand{\miusquarp}{{\color{petunidin!50}\scriptsize $\blacktriangleright$}~}
\newcommand{\niu}{\item[\miusquare]}
\newcommand{\nip}{\item[\miusquarp]}

%%%%% ABSAETZE

\newcommand{\pps}{\par \vspace*{1mm}}
\newcommand{\pp}{\par \vspace*{3mm}}
\newcommand{\ppbig}{\par \vspace*{8mm}}
\newcommand{\pphuge}{\par \vspace*{10mm}}

%================================
% FRAGE BOX
%================================

\tcbuselibrary{theorems,skins,hooks}
\newtcbtheorem{qst}{Frage:}
{%
	enhanced,
	breakable,
	colback = gray!15!white,
	frame hidden,
	boxrule = 0sp,
	borderline west = {2.5pt}{0pt}{red!70},
	sharp corners,
	detach title,
	before upper = \tcbtitle\par\smallskip,
	coltitle = black,
	fonttitle = \bfseries\sffamily,
	description font = \mdseries,
	separator sign none,
	segmentation style={solid, gray!50!black},
}
{th}




%================================
% ZITAT BOX
%================================

\tcbuselibrary{theorems,skins,hooks}
\newtcbtheorem{zitat}{Zitat}
{%
	enhanced,
	breakable,
	colback = gray!15!white,
	frame hidden,
	boxrule = 0sp,
	borderline west = {2.5pt}{0pt}{black!70},
	sharp corners,
	detach title,
%	before upper = \tcbtitle\par\smallskip,
	coltitle = gray!50!black,
	fonttitle = \bfseries\sffamily,
	description font = \mdseries,
%	separator sign none,
	segmentation style={solid, gray!50!black},
}
{th}



%================================
% SIMPLE SCHWARZER RAHMEN BOX
%================================


\newcommand{\boxfr}[1]{
	\begin{tcolorbox}[
	drop shadow southeast,
	enhanced,
	colframe=black!50,
	colback=white,
	boxrule = 1pt, 
%	toprule=0pt, bottomrule=0pt,rightrule=0pt,leftrule=0pt,
	boxsep=5pt,
	arc=1pt,
	]
	\begin{center}
	#1
	\end{center}
	\end{tcolorbox}
}

\definecolor{miudiutex}{HTML}{1d2926}
\definecolor{miugray}{HTML}{b4cccf}
\definecolor{miugray2}{HTML}{4d7365}
\definecolor{white2}{HTML}{f5f7f8}
\definecolor{red}{HTML}{cf1f5c}
\definecolor{green}{HTML}{00A36C}
\definecolor{delphinidin}{HTML}{315794}
\definecolor{uniblau}{HTML}{004e9f}
\definecolor{unigelb}{HTML}{fcba00}
\definecolor{unigrau}{HTML}{909085}
\definecolor{pelargonidin}{HTML}{b33807}
\definecolor{cyanidin}{HTML}{ba0746}
\definecolor{paeonidin}{HTML}{d15ca6}
\definecolor{petunidin}{HTML}{b56fed}
\definecolor{malvidin}{HTML}{8a2d8c}

\newcommand{\metermilli}{\mskip3mu \text{mm}}
\newcommand{\cm}{\mskip3mu \text{cm}}
\newcommand{\m}{ \text{m}}
\newcommand{\lite}{ \text{l}}
\newcommand{\kg}{ \text{kg}}
\newcommand{\g}{ \text{g}}
\newcommand{\km}{ \text{km}}
\newcommand{\h}{ \text{h}}
\newcommand{\s}{ \text{s}}
\newcommand{\tonne}{ \text{t}}
\usepackage[utf8]{inputenc}
\usepackage[T1]{fontenc}
\usepackage{lmodern}

\usepackage[
    activate={true,nocompatibility},
    final,
    tracking=true,
    kerning=true,
    spacing=true,
    factor=1100,
    stretch=30,
    shrink=20
]{microtype}
\usepackage{pdfpages}
\usepackage{awesomebox}
\usepackage{marginnote}
\tcbuselibrary{documentation}
\usepackage{sidenotes}
\usepackage{import}
\usepackage{xifthen}
\usepackage{transparent}
\usepackage[table]{xcolor}


\newcommand{\abbicon}{%
  \protect\tikz[baseline=0.2mm,scale=0.8]{%
    % Das blaue Quadrat mit abgerundeten Ecken
    \fill[uniblau, rounded corners=1.2pt] (0,0) rectangle (0.8em, 0.8em);
    % Der weiße Kreis in der Mitte
    \fill[white] (0.4em, 0.4em) circle (0.12em);
  }%
} 
\usepackage[labelfont={bf},font={color=miudiutex,small},figurename={\abbicon $~$Abbildung}]{caption}


\newcommand{\bckgrnd}{
\AddToHookNext{shipout/background}{
\begin{tikzpicture}[remember picture, overlay]
 \draw[draw=none,fill=uniblau!70, shift={(current page.north west)}] (-3,0) rectangle (23,-3.75);
 \draw[draw=none,fill=miugray, shift={(current page.north east)}] (-7,0) rectangle (3,-3.75);
\end{tikzpicture}
}
}



\newcommand{\incfig}[2]{%
\def\svgwidth{#2\columnwidth}
\import{C://LaTeX-Files/pngs/}{#1.pdf_tex}
}

\renewcommand{\textbf}[1]{\textsf{{\bfseries {#1}}}}



\titleformat{\section}[hang]
		{\LARGE \color{uniblau!75!black} \sffamily \bfseries}
		{\thesection}
		{6mm}
		{}
		[]
\titleformat{\subsection}[hang]
		{\large \sffamily \bfseries \color{black}}
		{\thesubsection}
		{4mm}
		{{\color{miugray}$\blacksquare ~$}}

\titleformat{\subsubsection}[block]
		{\normalsize \bfseries \sffamily \color{miudiutex}}%
		{\normalsize \color{black} \thesubsubsection}
		{2mm}
		{{\color{miugray}$\blacksquare ~$}}
		[ \color{black}]
		
\titleformat{\paragraph}
		{\normalsize \bfseries \sffamily \color{black}}
		{\footnotesize \theparagraph}
		{1mm}
		{}

\pagestyle{empty}
\titlespacing*{\section}{0pt}{20mm}{12mm}
\titlespacing{\section}{0pt}{20mm}{12mm}

\titlespacing*{\subsection}{0pt}{14mm}{4mm}
\titlespacing{\subsection}{0pt}{14mm}{4mm}

\titlespacing*{\subsubsection}{0pt}{6mm}{3mm}
\titlespacing{\subsubsection}{0pt}{6mm}{3mm}

\titlespacing*{\paragraph}{0pt}{6mm}{2mm}
\titlespacing{\paragraph}{0pt}{6mm}{2mm}


\let\oldmarginfigure\marginfigure
\let\endoldmarginfigure\endmarginfigure

\let\oldmarginfigure\marginfigure
\let\endoldmarginfigure\endmarginfigure

\renewenvironment{marginfigure}[1][0pt] % [1] steht für 1 Argument, [0pt] ist der Standardwert
  {\oldmarginfigure[#1]\captionsetup{labelfont={sf,bf,color=uniblau!75!black},font={color=miudiutex,small}}}
  {\endoldmarginfigure}
  
  
  
  
\renewcommand{\labelitemi}{\color{uniblau!70}$\bullet$}
\renewcommand{\labelitemii}{\color{uniblau!70}$\cdot$}
\newcommand{\plmtsquare}{{\color{uniblau!70}$\blacksquare~$}}
\newcommand{\splmt}[1]{{\color{uniblau!75!black} \sffamily #1}}
\newcommand{\fplmt}[1]{{\color{uniblau!75!black} \sffamily \bfseries #1}}

\newcommand{\antwort}[1]{
\begin{tcolorbox}[
	enhanced,
	enforce breakable,
	standard jigsaw,
	boxrule=0.7pt,
	colframe=black,
	sharp corners,
	boxsep=5pt,
	title=\bfseries \sffamily Antwort,
	opacityback=0,
	coltitle=miudiutex,
	attach boxed title to top text left={yshift=-\tcboxedtitleheight/2},
	boxed title style={colback=white,colframe=white},
	bottomrule at break=0pt,
	toprule at break=0pt,
	pad at break=16mm,
]
	\begin{center}
	#1
	\end{center}
	\end{tcolorbox}
}




\rowcolors{2}{miugray!50}{miugray!50}
\arrayrulecolor{white} % Weiße Gitterlinien
\setlength{\arrayrulewidth}{1.5pt} % Etwas dickere Linien
\renewcommand{\arraystretch}{1.4} % Mehr Zeilenhöhe für bessere Lesbarkeit}


%\setlist[itemize]{itemsep=0mm}
\setlist[enumerate]{itemsep=0mm}
\usepackage[table]{xcolor}
\usepackage{pgfplots}
    \pgfplotsset{
        every axis post/.style={
	yticklabel=\empty,
	xticklabel=\empty,
	y label style={at={(axis description cs:-0.2,1.1)}},
	x label style={at={(axis description cs:1.05,0.2)}},
	ticks=none,
	width=6cm,
	axis lines=middle,
%	xlabel near ticks,
        },
    }
\usetikzlibrary{positioning, arrows.meta, calc}

\usepackage{relsize}
\usepackage{cancel}

\newif\ifshowme
%\showmetrue



\begin{document}

%\newgeometry{top=3cm,bottom=3cm,left=4.5cm,right=4.5cm}
\newgeometry{top=2cm,bottom=2cm,left=4cm,right=4cm}


\begin{center}
\LARGE
\color{uniblau}
\textbf{Übung 06} \par \vspace*{2mm}
\large \color{miugray!50!black} \textbf{Prozessbezogene Lebensmitteltechnologie}\par 
\normalsize Zusatzaufgaben
\end{center}
\par 
%\vspace*{6mm}


\color{miudiutex}

%\begin{multicols}{2}
\subsection*{Aufgabe 1}
Was ist der Unterschied zwischen Reindichte, Rohdichte und Schüttdichte?


\subsection*{Aufgabe 2}
In \figurename{~\ref{masse}} sehen Sie den fünfstufigen Herstellungsprozess eines Fertigprodukts. 
\begin{itemize}
\item Bestimmen Sie den Massenstrom von Produkt $P$
\item Bestimmen Sie den Massenstrom von Rücklauf $A$
\end{itemize}


\begin{figure}[h]
\centering
\resizebox{.8\textwidth}{!}{
\begin{tikzpicture}[font=\sffamily, >={Latex[length=6pt, width=5pt]}]

    % --- Definition der Stile für die Blöcke ---
    \tikzset{
        block/.style={rectangle, draw=black, minimum width=1.6cm, minimum height=1.4cm, align=center},
		blaublock/.style={block, fill=uniblau!50},
        grayblock/.style={block, fill=miugray!50},
        redblock/.style={block, fill=miugray}
    }

    % --- Positionierung der Blöcke ---
    \node[blaublock] (I)   at (0, 0)      {I};
    \node[redblock]  (II)  at (0, -2.5)   {II};
    \node[redblock]  (III) at (-1.5, -5.0){III};
    \node[redblock]  (IV)  at (1.5, -5.0) {IV};
    \node[grayblock] (V)   at (1.5, -7.5) {V};

    % --- Systemgrenze (gestrichelte Linie) ---
    \draw[dashed, semithick, uniblau] (-3.2, 1) rectangle (3.2, -8.7);
    \node[above left, inner sep=2pt] at (3.2, 1.2) {\color{uniblau}Systemgrenze};

    % --- Stoffströme (Pfeile & Beschriftungen) ---
    
    % Feed F
    \draw[->] (-4.8, 0) -- (-2.7, 0);
    \draw[->] (-2.7, 0) -- (I.west);
    \node[left, align=left, inner sep=2pt] at (-5, 0) {Feed F \\ 1000 kg/h\\15\% TS\\85\% Wasser};

    % Stream D
    \draw[->] (I.east) -- (4.2, 0);
    \node[right, align=left, inner sep=2pt] at (4.5, 0) {D \\ 150 kg/h\\0\% TS};

    % Stream B
    \draw[->] (I.south) -- (II.north) node[midway, right] {B};

    % Stream C, E und G (Aufteilung)
    \draw[->]  (II.south) -- (0, -3.75) node[midway, right] {C};
    \draw[-]  (0, -3.75) -- (-1.5, -3.75) node[midway, below,xshift=5pt] {E 10\% TS};
    \draw[->] (-1.5, -3.75) -- (III.north);
    \draw[-]  (0, -3.75) -- (1.5, -3.75) node[right] {G};
    \draw[->] (1.5, -3.75) -- (IV.north);

    % Stream A (Rückführung)
    \draw[-]  (III.west) -- (-2.7, -5.0);
    \draw[->] (-2.7, -5.0) -- (-2.7, 0);
    \node[right, align=left] at (-2.7, -2.5) {A\\5\% TS};

    % Stream R (Rückführung)
    \draw[-]  (IV.east) -- (2.7, -5.0) -- (2.7, -2.5);
    \draw[->] (2.7, -2.5) -- (II.east);
    \node[left] at (2.7, -3.75) {R};

    % Stream L
    \draw[->] (IV.south) -- (V.north) node[midway, right] {L};

    % Stream K
    \draw[->] (III.south) -- (-1.5, -9.0) node[below, align=left] {K\\450 kg/h\\20\% TS};

    % Stream Product P
    \draw[->] (V.south) -- (1.5, -9.0) node[below, align=left] {Produkt P\\80\% TS};

    % Stream W
    \draw[->] (V.east) -- (4.2, -7.5) node[right, align=left] {W\\Wasser};

\end{tikzpicture}
}

\caption{Schematische Darstellung des Herstellungsprozesses. Abgeändert nach Singh und Heldman.}
\label{masse}
\end{figure}

\pagebreak
\subsection*{Aufgabe 3}
Sie veranstalten eine Sommerparty und Daniel Düsentrieb schenkt Ihnen eine absolut perfekt isolierte Kühlbox. Leider hat er vergessen auch den Styropor-Deckel perfekt zu isolieren.
Der Deckel ist 0,5 m$^2$ groß und besitzt eine Dicke von 10 cm. 
Die Getränke in der Kühlbox werden auf 7° C gemessen. Die Temperatur außerhalb der Kühlbox beträgt 35 °C. Die Wärmeleitfähigkeit des Styropors ist $\lambda_1$ = 0,04 W/mK. Auf der Innenseite des Deckels ist eine zusätzliche 2 cm dicke Isolierschicht mit Wärmeleitfähigkeit $\lambda_2$ = 0,02 W/mK angebracht. 

\begin{itemize}
\item Wie lange dauert es bis die 10 Stck. 0,33 L Flaschen Spaßgetränke (Wärmekapazität c = 4,2 J/kg K) zu warm sind (15 °C) ?
\item Welche Temperatur herrscht an der Grenze zwischen Isolierschicht und Styropor?
\end{itemize}

\end{document}
